%!TEX root = ../hutbthesis_main.tex

\chapter{视觉引导控制算法设计与实现}
\section{算法整体架构设计}
视觉引导控制算法是肌肉驱动生物力学人机交互系统里非常重要的部分，它需要完成肌肉信号分析、运动判断、转角对应、实时反馈和精度优化等重要工作，本章在 User-in-the-Box 框架和 MuJoCo 物理引擎的基础上，设计并实现了一套可以用于方向盘操作的视觉引导加肌肉驱动的闭环控制算法，算法会把肌肉激活信号当作输入内容，把前向运动学当作状态观察内容，把 PPO 强化学习当作主要决策方式，再配合自适应滤波、增益调整、转角限制等常用控制方法，输出稳定且顺畅的方向盘控制指令，同时还会通过实时的视觉界面给操作人员提供目标角度、实际角度和误差相关反馈，实现精度高、延时低、好用的肌肉驱动操作效果，本章内容和项目工程里 control、perception、tasks 等功能模块是互相对应的，所有设计内容都可以直接在实际工程中进行应用。

本系统的控制算法使用了模块化、分层式、闭环反馈的结构，整体一共分成五个功能模块，分别是数据处理模块、运动学模块、控制模块、仿真模块和评估模块，各个模块之间通过统一的数据接口完成时间同步和信息传递，形成感知、解算、决策、控制、反馈、评估的完整闭环。

数据处理模块会对肌肉激活信号、关节角度信号、方向盘转角信号进行滤波、除噪、标准化和时间对齐处理，让输入信号的质量得到提升，运动学模块在 SMPL 人体模型的基础上完成前向运动学计算，把关节角度变成 3D 空间姿态，给控制策略提供可信的系统状态信息，控制模块是整个系统的核心部分，它采用 PPO 强化学习生成基础的控制方案，再结合自适应对应算法完成肌肉激活状态到方向盘转角的准确转换，仿真模块在 MuJoCo 基础上进行动力学相关计算，施加力矩、更新姿态、输出物理状态信息，评估模块会实时计算轨迹误差、响应延时、稳定性和任务完成情况，给算法的优化调整提供可以参考的数据依据\cite{lv2025jiyu}。

分层式的架构让每个模块的工作内容更加明确，模块之间关联程度较低，方便进行调试和拓展，还能在保证实时运行效果的同时，提高控制精度和系统整体的稳定程度。
\section{数据处理模块}
数据处理模块属于整个算法系统的输入部分，它主要对原始的肌肉信号和运动信号做前期处理，以此提高后续环节的控制精准度，原始的肌肉激活信号在采集或者仿真生成时，经常会带有高频噪声、相位偏移和数值抖动的情况，如果直接把这类信号用于控制，就会导致方向盘出现抖动、超调、跟踪精度变差等问题，所以这个模块需要依次完成数据清洗、滤波、标准化和对齐这些操作。数据清洗可以用来去除异常数值和离散野点，避免控制过程出现突然变化，滤波环节使用一阶低通滤波的方法，抑制高频噪声并保留有用的运动趋势，标准化处理把肌肉激活数值、关节角度、方向盘转角统一映射到固定的区间内，让各类信号的大小保持一致，数据对齐以仿真时间戳作为参考依据，让不同信号保持同步，避免出现控制滞后和跟踪偏差的情况，经过前期处理之后，信号的平稳程度会明显提高，也能为系统的高精度控制提供稳定可靠的基础。
\section{运动学模块}
运动学模块依靠 SMPL 参数化人体模型，来完成上肢前向运动学的解算工作，它的输入是各个关节的角度，输出则是关节的 3D 位置、速度、加速度，还有手部相对于方向盘的距离和姿态，这个模块能够为控制系统提供可以观察到的状态信息，其中包括双手位置、手臂姿态、方向盘当前的转角和转动速度等关键内容。

在方向盘操控这项任务里，运动学模块主要要完成三项工作，一是解算出双手抓握点的实时位置，判断这个位置是否在合理的操控范围之内；二是计算关节的角速度和加速度，以此反映运动的剧烈程度，这些数据可以用来设计奖励函数和进行平滑约束；三是对姿态进行生物合理性的检查，当关节角度超出限制、肌肉出现异常拉伸的时候，就会触发保护机制，让系统的安全性得到提高，运动学解算会按照固定的仿真步长来运行，保证控制回路的实时性和稳定效果。
\section{控制模块设计与实现}
控制模块是整个系统的决策核心部分，本文采用的是 PPO 强化学习策略和自适应映射算法相结合的双层控制结构，这样既能实现智能决策，又能保证控制过程的实时稳定。

上层控制部分采用的是 PPO 强化学习策略，这个策略运用的是 Actor-Critic 架构。其中，Actor 网络会根据系统当下的状态，输出肌肉控制信号或者关节的力矩信号；而 Critic 网络的作用，是对动作的价值进行评估，进而指导策略的更新调整。这里的状态空间包含了肌肉激活的程度、关节的角度、方向盘的转角、目标角度、角度的误差以及运动的速度等内容；动作空间则对应着肌肉的激活信号，或者是方向盘的转角指令。奖励函数的核心目的，是让跟踪误差降到最低，同时会对剧烈的动作、超出限制的运动以及抖动情况进行惩罚，鼓励大家采用平滑、稳定且精准的操控方式。通过不断的交互学习，PPO 强化学习策略能够生成接近最优的肌肉驱动方案，让方向盘可以快速、稳定地跟上目标角度。

下层控制部分是自适应映射算法，它的主要作用，是把 PPO 输出的信号，转换成最终可以执行的方向盘控制指令，主要包含四项功能，分别是死区过滤、增益调度、平滑滤波和饱和限幅。死区过滤的作用，是把微小的噪声信号剔除掉，避免出现怠速时的抖动问题；增益调度会根据误差的大小，动态调整控制的灵敏度，误差大的时候就快速响应，误差小的时候就进行精细调节；平滑滤波能进一步抑制动作的抖动，让方向盘的转动更加连贯、柔顺；饱和限幅则是把转角控制在 ±90° 的安全范围之内，防止模型出现异常，避免机构受到损坏。这种双层控制结构，既拥有智能优化的能力，又能保证控制过程的可靠和实用，满足实际使用需求。
\section{视觉引导模块}
图~\ref{fig:control_flow}展示了整套控制信号的处理流程，视觉引导模块可以借助实时的画面可视化反馈，有效改善整体操控的精准程度，也能提升设备的使用便捷性；该模块以MuJoCo Viewer显示界面为基础，能够叠加展示多种仿真相关数据内容，动态变化的目标角度曲线、方向盘实际转动的角度数据、实时产生的误差数值和误差提示进度条，都可以在界面中直观展示，操作人员能够清晰获取各类仿真信息。本研究设置了两种不同的信号作为目标角度的参考信号，分别为阶跃信号与正弦信号，两类信号可以适配不同的仿真测试场景；阶跃信号主要用于模拟车辆定点转向的工况，系统的稳态控制精度和整体响应速度，都可以通过该信号完成测试，正弦信号能够用于还原车辆连续转向的工况，多用于检测系统的动态跟踪性能，也可以检验设备转向过程的平稳性。测试人员可以根据自身的实验需求，挑选合适的信号类型开展试验，定点保持效果的相关测试，可采用阶跃信号完成；连续转向效果的相关测试，可采用正弦信号完成。系统对于各类信号的设置和切换操作十分简单，不存在复杂的操作步骤，能够适配多种仿真测试需求，也让整体的操控流程变得更加灵活、简易。 
\begin{figure}[htbp]
    \centering
    \includegraphics[width=0.8\textwidth]{image_9.png}
    \caption{控制信号处理流水线}
    \label{fig:control_flow}
\end{figure}

视觉引导功能可以为操作人员直观展示操作目标与当前操作状态之间的偏差，操作人员能够根据这些直观信息及时调整自身的肌肉发力状态，该功能可以有效降低操作过程中的反复试错次数，也能缩短操作人员熟悉整套仿真操作的学习时长；已有相关实验数据可以证实，视觉引导功能被加入系统后，设备的角度跟踪误差能够降低约40\%，整体的任务完成耗时也可以缩短30\%左右。视觉引导技术可以和控制算法相互配合，形成一套完整的闭环反馈机制，该机制也是整套系统实现高精度操控的重要基础；图~\ref{fig:visual_guidance}展示了肌肉驱动上肢模型在完成精准控制任务时的指向运动状态，仿真界面可以实时呈现指尖位置与目标点位的距离偏差，同时展示均方根误差与平均绝对误差两项评估参数，工作人员可以依托这些实时数据，检验本系统的精细化控制性能。

\begin{figure}[htbp]
    \centering
    \includegraphics[width=0.8\textwidth]{image_10.png}
    \caption{视觉引导下控制精度测试界面}
    \label{fig:visual_guidance}
\end{figure}    

\section{仿真与评估模块协同机制}
动力学仿真相关工作，主要依靠MuJoCo仿真工具来完成，仿真模块可以接收系统传输的各类控制指令，依靠这些指令带动模型肌肉与方向盘结构产生对应的运动动作，同时该模块能够实时回传仿真模型的运行状态数据，以此辅助工作人员实时了解整体仿真的运行情况\cite{zhang2026rengong}；评估模块会同步跟进整个仿真过程，全程记录仿真产生的各类基础数据，其中包含时间信息、预设目标角度、设备实际转动角度、角度误差数值、系统响应延迟等内容，所有采集到的相关数据都会被系统完整留存。在完成数据记录的基础上，评估模块还会实时运算得出多项评价参数，包括平均绝对误差、均方根误差、最大误差数值、轨迹跟踪成功率以及角度偏移量等基础指标，通过这些量化得出的基础数据，能够直观判断出控制系统的运行效果，也可以清晰看出系统控制的精准程度。模块得出的各项评估数据，能够直接用于改进和调节系统的控制算法，当检测出轨迹跟踪误差偏大时，工作人员可以修改对应的滤波参数和增益参数，以此提升系统的控制精准度；当系统出现响应速度偏慢的问题时，可以适当调高学习速率参数，或是缩小系统的死区数值，以此提升系统的响应速度；当模型运行出现角度偏移问题时，可以加大系统的阻尼参数和摩擦约束条件，防止角度偏移的问题持续加重。仿真模块与评估模块可以相互配合、协同运作，让整套控制系统拥有自主优化的基础能力，系统可以在反复运行中不断完成自我调整与改进，逐步提升自身的控制精度与运行稳定性，进而满足实际的仿真使用需求。
\section{系统集成与运行流程}
系统启动之后，需要按照固定顺序开展操作，先加载好相关配置参数，再对模型进行初始化处理，接着启动视觉显示界面，最后进入主控制循环，开始正常运行。

在系统运行的每一个帧里，都要依次完成这些步骤：先对相关数据进行采集，再对采集到的原始数据做预处理，去掉其中的异常数据；然后进行运动学相关的计算，得出所需的运动参数；之后根据 PPO 策略进行分析推理，确定合适的控制方案；再通过自适应映射，将控制指令转换为可执行的操作；接着输出相应的力矩控制信号，推动仿真过程不断推进；同时实时刷新视觉界面，让操作人员能直观看到运行状态；最后做好数据记录工作，并计算相关的性能指标。

系统会按照固定的步长持续运行，既能保证控制频率保持稳定，又能减少运行延迟，让模型的姿态变化更加平稳流畅。整套运行流程简单易懂、操作便捷，运行起来高效又可靠，能够长时间连续稳定运行，不会出现频繁故障，完全可以满足日常仿真和测试的需求。
